\documentclass[11pt,notitlepage]{report}

\usepackage[T1]{fontenc}
\usepackage{lmodern}
\usepackage[margin=1in]{geometry}
\usepackage{amsmath}
\usepackage{amssymb}
\usepackage{amsthm}
\usepackage{microtype}
\usepackage[colorlinks=true, allcolors=blue]{hyperref}
\usepackage{xurl}
\setlength{\emergencystretch}{2em}

\theoremstyle{definition}
\newtheorem{definition}{Definition}
\newtheorem{algorithm}{Algorithm}

\theoremstyle{plain}
\newtheorem{lemma}{Lemma}

\setlength{\parindent}{0pt}
\setlength{\parskip}{0.5em}

\title{NN Architecture notes}
\author{}
\date{}

\begin{document}

\maketitle

\tableofcontents

\chapter{Stanford CS231n}
\url{https://cs231n.github.io/neural-networks-1/#neural-network-architectures}

I already know a lot of this so going to skip most of it. Weights and biases.
They say that for classic Feed-Forward (they call this MLP right?) that 3 layers
is usually "good enough" and having more is not often helpful. This strikes me
as absolutely wild. But I've never actually built MLP before so who am I to say.

They mention a few ways to prevent overfitting but don't explain them.
(L2 regularization, dropout, input noise). I at least sort-of know what those are
anyways so it's not too big a deal.

Apparently some research has shown that smaller networks are more likely to get
stuck in a bad local minimum. Very interesting, I wonder why.

Overall I learned ~nothing from this page.

\chapter{Pandas Preprocessing}
\url{https://d2l.ai/chapter_preliminaries/pandas.html}
WOW! WAY TOO low-level. It says to convert your features into a vector of real numbers.
That's a really useless page.

\chapter{More Stanford CS231n}
\url{https://cs231n.github.io/neural-networks-2/}

This is interesting background but not useful for my projects.
They're talking about various preprocessing, from simple to complicated.
Zero-center, normalization are the simple.
They then talk about PCA, which feels a bit crazy that my PhD shows up here.
It's been so long though that I don't remember everything. Something something
grab some eigenvectors of somebody.
Reduce the dimensionality by only using the top eigenvalues. Nice.
The also suggest to divide by eigenvalues along the eigenbasis as a way of
REALLY REALLY normalizing.

\section{Weight Initialization}

Okay this actually matters for my project. How to initialize the weights?
\begin{itemize}
\item All-zero initialization = really stupid. Once you have symmetry you can't go back
\item Small random numbers = good start. Unclear if they should be small or not. (Small number = small gradient for backprop)
\item Small random numbers + normalized variance.
If we just have small random numbers then variance of a neuron increases with fan-in.
So depending on layer dimensions... I guess your training is non-uniform? Hard to imagine
EXACTLY the issue but it makes sense this could be bad. If one layer gets too little update
and another layer gets too much.

Instead divide by $1/\sqrt{n}$, where $n$ is the fan-in. Somebody else scale variance with
$2/\left(n_{in} + n_{out}\right)$. I think this is getting over my head for now. Good to know
the starting details and I'll learn more later (maybe).
\item Sparse initialization. Set almost all weights to zero, but each node has non-zero (random)
initialized connection to ~10 nodes below it. So it's not symmetric anymore but still most stuff is zero.
\item Bias initialization. This page seems unsure what to do. Some recommendations are
to set it all to zero, or to set it all to some low non-zero number (so all the ReLU fire)
\item Batch Normalization. Add a BatchNorm layer immediately after each fully connected layer,
before nonlinearities. Since BatchNorm is differentiable this works great for your gradient
descent. TODO: learn the details of BatchNorm (I have a guess but I'm sure it's not exact).
\end{itemize}

\section{Regularization}

Seems I need to learn some more vocab before I can understand this.
Let's quit here and come back later.

\begin{itemize}
\item L2 Regularization
\item L1 Regularization
\item Max norm constraint
\item Dropout
\end{itemize}

\chapter{Gymnasium}
\url{https://gymnasium.farama.org/main/introduction/create_custom_env/?utm_source=openai}

\end{document}
